\section{Problem \#654}

\subsection*{1. ROUND-2 OBJECTIVE}

\textbf{Path (B): explicit counterexample / disproof.}

Round 1 left the main gap completely open: no counterexample family was known there.  In Round 2 I give an explicit infinite family showing that the strongest form of the conjecture is false.

\subsection*{2. Round-1 FOUNDATION USED}

I use only the Round-1 definition of
\[
f(n):=\min_{|X|=n}\max_{x\in X} R_X(x)
\]
under the constraint that $X\subset \mathbf{R}^2$ has no four points on a circle, together with the trivial lower bound $f(n)\ge (n-1)/3$.

No nontrivial Round-1 lemma is needed for the disproof.

\subsection*{3. NEW INSIGHT / TOOL (ROUND 2)}

The new construction has two ingredients.

\begin{itemize}
\item Put symmetric pairs of points on each coordinate axis.  Then from any point on one axis, the distances to the opposite axis come in equal pairs, so the cross-axis contribution is halved.
\item Choose the $x$-coordinates to be powers of $2$ and the $y$-coordinates to be powers of $3$.  Then any product of two signed $x$-coordinates is a signed power of $2$, while any product of two signed $y$-coordinates is a signed power of $3$.  This forbids the only possible $4$-point concyclicity pattern.
\end{itemize}

\subsection*{4. ATTACK PLAN (ROUND 2)}

For each $t\ge 1$, define
\[
X_t:=\{(\pm 2^i,0):1\le i\le t\}\cup\{(0,\pm 3^j):1\le j\le t\}.
\]
Then $|X_t|=4t$.

I prove two claims.

\begin{enumerate}
\item No four points of $X_t$ lie on a common circle.
\item Every point of $X_t$ determines at most $3t-1$ distinct distances to the other points.
\end{enumerate}

These two claims imply
\[
f(4t)\le 3t-1=\frac{3}{4}(4t)-1,
\]
which disproves the asymptotic statement $f(n)>(1-o(1))n$.

\subsection*{5. WORK (ROUND 2)}

Fix $t\ge 1$, and define
\[
A_t:=\{\pm 2^i:1\le i\le t\},\qquad B_t:=\{\pm 3^j:1\le j\le t\},
\]
\[
X_t:=\{(a,0):a\in A_t\}\cup \{(0,b):b\in B_t\}.
\]
Then $|X_t|=4t$.

\paragraph{Claim 1: no four points of $X_t$ are concyclic.}
A circle meets any fixed line in at most two points.  Since all points of $X_t$ lie on the two coordinate axes, any set of four concyclic points from $X_t$ would have to consist of exactly two points on the $x$-axis and exactly two points on the $y$-axis.  So it would have the form
\[
(a,0),(a',0),(0,b),(0,b')
\]
with $a\ne a'$ and $b\ne b'$.

Suppose such four points lie on a circle.  Write the equation of the circle as
\[
x^2+y^2+ux+vy+w=0.
\]
Plugging in $(a,0)$ and $(a',0)$ shows that $a$ and $a'$ are roots of
\[
x^2+ux+w=0,
\]
so by Vi\`ete's formula,
\[
aa'=w.
\]
Similarly, plugging in $(0,b)$ and $(0,b')$ shows that $b$ and $b'$ are roots of
\[
y^2+vy+w=0,
\]
so
\[
bb'=w.
\]
Hence
\[
aa'=bb'.
\]
But $aa'$ is a signed power of $2$, while $bb'$ is a signed power of $3$.  These can never be equal.  Contradiction.

Therefore no four points of $X_t$ lie on a common circle.

\paragraph{Claim 2: every point of $X_t$ sees at most $3t-1$ distinct distances.}
Take first a point $p=(a,0)$ on the $x$-axis.

\begin{itemize}
\item There are $2t-1$ other points of $X_t$ on the $x$-axis, so their contribution to $R_{X_t}(p)$ is at most $2t-1$.
\item For each $j\in\{1,\dots,t\}$, the two points $(0,3^j)$ and $(0,-3^j)$ are both at distance
\[
\sqrt{a^2+3^{2j}}
\]
from $p$.  So all $2t$ points on the $y$-axis contribute at most $t$ distinct distances.
\end{itemize}
Hence
\[
R_{X_t}(p)\le (2t-1)+t=3t-1.
\]

The same argument applies to a point on the $y$-axis by symmetry.  Thus every point of $X_t$ determines at most $3t-1$ distinct distances.

So for $n=4t$ we have constructed an $n$-point set with no four points on a circle and with
\[
\max_{x\in X_t}R_{X_t}(x)\le 3t-1=\frac{3n}{4}-1.
\]
Therefore
\[
f(4t)\le \frac{3}{4}(4t)-1.
\]
In particular,
\[
\limsup_{t\to\infty}\frac{f(4t)}{4t}\le \frac{3}{4}<1.
\]
So the statement $f(n)>(1-o(1))n$ is false.

\subsection*{6. ADVERSARIAL VERIFICATION}

\begin{itemize}
\item The concyclicity proof covers all sign patterns: any product of two signed $x$-coordinates is $\pm 2^m$, and any product of two signed $y$-coordinates is $\pm 3^n$.
\item A circle cannot contain three points from one axis and one from the other, because each axis meets a circle in at most two points.
\item In the distance count I used only upper bounds, so possible accidental overlaps between same-axis and cross-axis distances can only help.
\item The family is given for $n=4t$.  That is enough to disprove an asymptotic assertion of the form $(1-o(1))n$, since such a statement must hold along every infinite subsequence.
\end{itemize}

\subsection*{7. FINAL}

\textbf{FULL SOLUTION -- COUNTEREXAMPLE / DISPROOF}

The strongest form of Problem \#654 is false.  The explicit family
\[
X_t=\{(\pm 2^i,0):1\le i\le t\}\cup\{(0,\pm 3^j):1\le j\le t\}
\]
has no four points on a circle and satisfies
\[
\max_{x\in X_t}R_{X_t}(x)\le \frac{3}{4}|X_t|-1.
\]
Thus $f(n)$ is not $(1-o(1))n$ in general.  This agrees with the current status note on the Erd\H{o}s Problems site \cite{EP654}.

\subsection*{8. COMPLETION ESTIMATE}

\textbf{COMPLETION: 100\%}
